\documentclass[10pt,a4paper]{amsart}
\usepackage[margin=19mm]{geometry}
\usepackage{amsmath,amssymb,mathtools}
\usepackage[scr=rsfs,cal=euler]{mathalfa}
\usepackage{xfrac}
\usepackage[unicode,hidelinks,bookmarks=false]{hyperref}
\newcommand{\ie}{i.e.\ }
\newcommand{\cf}{cf.\ }
\newcommand{\resp}{respectively\ }
\newcommand{\Subsection}{\subsection}
\providecommand{\nobreakdash}{\nobreak\hbox{-}}
\setlength{\emergencystretch}{2em}
\multlinegap0pt
\newtheorem{theorem}{Theorem}
\def \s{{\mathbb S}}
\title{Conformal-biharmonic hypersurfaces in spheres and product spaces}
\author{V. Branding, S. Montaldo, S. Nistor, C. Oniciuc, A. Ratto}
\date{Source: arXiv:2502.02510v2 (CC BY 4.0)}
\begin{document}
\maketitle

\noindent\emph{Source and licence.} Conformal-biharmonic hypersurfaces in spheres and product spaces. Version arXiv:2502.02510v2 (CC BY 4.0), distributed by arXiv under the Creative Commons Attribution 4.0 International licence, http://creativecommons.org/licenses/by/4.0/, as recorded in the arXiv licence metadata for this exact version and shown on the abstract page https://arxiv.org/abs/2502.02510v2. Full licence notice: \texttt{licenses/arxiv-2502.02510-CC-BY-4.0.txt}.\par\smallskip

\noindent\textit{Editorial scope.} Counted unit: the article's theorem Th-c-biharm-not-minimal-l=4 (source0.tex lines 549--551). The article's complete original proof of it is delivered with it (source0.tex lines 552--614, one \texttt{proof} environment of 63 lines). No mathematics is written, repaired or re-derived: the statement and the proof are the article's own lines, in the article's own notation, and no retained line is substituted or re-flowed.  No further source-local context is needed: every cross-reference of every retained line resolves inside the two delivered spans.  Nothing was omitted from any retained span, so the package carries no omission record.  No retained line contains a citation, so the wrapper carries no bibliography and no bibliography entry.  No retained line contains a figure environment, an image-inclusion command or drawing code, so no image is delivered and the package declares no asset dependency.  Editorial formatting only, in addition: an AMS \texttt{amsart} preamble that restates the article's own declarations and macros used by the retained lines, namely the environments \texttt{theorem} on separate counters, together with the standard packages those lines need.  The retained environments print in this document as Theorem 1 in order of appearance, whereas the article prints the same items with the numbering of its own full document.  The complete native source of the article as distributed in the arXiv e-print for this exact version is bundled unchanged as \texttt{sources/172737/source0.tex}.
\begin{theorem}\label{Th-c-biharm-not-minimal-l=4}
Let $M_s^m$, $s\in(0,\pi/4)$, be a family of isoparametric hypersurfaces of degree $4$ in $\s^{m+1}$ with $m_1<m_2$, $m=2(m_1+m_2)$. Then there exists a unique $s^*\in(0,\pi/4)$ such that $M_{s^*}^m$ is an isoparametric $c$-biharmonic hypersurface in $\s^{m+1}$. Moreover, the $c$-biharmonic hypersurface $M_{s^*}^m$ is non-minimal.
\end{theorem}

\begin{proof}
Direct substitution and simplification yields:
\[
H=\frac{ \Big((m_1+m_2) \cos (4s)+m_1-m_2 \Big)}{\sin(4s) (m_1+m_2)}\,
\]
\[
[\tau_2^c]^\perp=\frac{1}{32 \sin^3(s) \cos^9(s)\left(\tan ^2(s)-1\right)^3} \,A(s)\,,
\]
where
\begin{equation}\label{eq:polyAs}
\begin{aligned}
A(s)=&2 (m_1 - m_2) \big(12 (m_1^2 + m_2^2) + 8 m_1 m_2 - 33 ( m_1 + m_2) + 36\big)\\
   &+ \big(32 (m_1^3 + m_2^3) - 83 (m_1^2 + m_2^2) - 6 m_1 m_2 + 96 (m_1 + m_2)\big)\cos (4 s)\\
   &+2 (m_1 - m_2) \big(4 (m_1^2 + m_2^2) + 8 m_1 m_2 - 7 (m_1 + m_2) + 12\big) \cos (8 s)\\
   &+3 (m_1 + m_2)^2 \cos (12 s)\,.
\end{aligned}
\end{equation}

Now, it is easy to check that there exists a unique value $\tilde{s}\in (0,\pi/4)$ such that $H(\tilde{s})=0$, i.e.,
\[
\tilde{s}= \frac{1}{4} \cos
   ^{-1}\left(\frac{m_2-m_1}{m_1+m_2}\right)\,.
\]
Next, we compute
\begin{equation}\label{eq:limiti}
\begin{aligned}
\lim_{s\to 0^+} A(s) =& 32 m_1 (6 + m_1 (-5 + 2 m_1))>0\\
\lim_{s\to \pi/4^-} A (s)  =&-32 m_2 (6 + m_2 (-5 + 2 m_2))<0\,.
\end{aligned}
\end{equation}
From this it is easy to deduce that $\lim_{s \to 0^+}[\tau_2^c]^\perp(s)=-\infty$ and $\lim_{s \to \pi/4 ^-}[\tau_2^c]^\perp(s)=+\infty$. 

Therefore, there exists $s^*\in (0,\pi/4)$ such that $[\tau_2^c]^\perp(s^*)=0$. 

Finally, 
\[
[\tau_2^c]^\perp(\tilde{s})= \frac{48 (m_1^2-m_2^2)}{\sqrt{m_1 m_2}}   \neq 0
\]
which implies $s^*\neq \tilde{s}$. Therefore $M_{s^*}^m$ is $c$-biharmonic and non-minimal.
To end the proof it remains to show that $s^*$ is the unique solution of $A(s)=0$ in $(0,\pi/4)$. We rewrite $A(s)$ in \eqref{eq:polyAs} as a polynomial by setting $y=\cos(4s)$.  We obtain
\begin{equation}\label{eq:polyAy}
A(y)=a_o+a_1y+a_2y^2+a_3y^3\,,
\end{equation}
where
\[
\begin{aligned}
a_0 &= 4 (m_1 - m_2) \big(4(m_1^2+m_2^2)-13(m_1+m_2)+12\big)\\
a_1 &=  4 \big(8 (m_1^3+ m_2^3) -23(m_1^2+m_2^2)-6m_1 m_2+24(m_1+m_2)\big)\\
a_2 &= 4 (m_1 - m_2) \big(4(m_1^2+m_2^2)+8m_1 m_2-7(m_1+m_2)+12\big)\\
a_3 &= 12 (m_1 + m_2)^2\,.
\end{aligned}
\]
We have to show that the polynomial $A(y)$ admits only one root in $(-1,1)$. We deduce from \eqref{eq:limiti} that  $A(-1)<0$  and $A(1)>0$.
Then, it is enough to show that $A'(y)$ admits at most one zero in $(-1,1)$. It is easy to check that $A''(-1)<0$ and that $A''(y)$ is strictly increasing. Next a straightforward analysis shows that
\[
\begin{cases}  \text{(i)} \quad A''(1)<0 \quad \text{if}\; m_2=m_1+k, \; k\geq 3\\ 
\text{(ii)} \quad A''(1)>0 \quad \text{if}\; m_2=m_1+1 \;\text{or}\; m_2=m_1+2.
\end{cases}
\]
In the case (i), we have $A''(y)<0$ for all $y\in (-1,1)$ and so $A'(y)$ is strictly decreasing on $(-1,1)$. Therefore, $A'(y)$ admits at most one zero in $(-1,1)$.

In the case (ii), there exists a unique point $y_0=-a_2/(3a_3)\in(-1,1)$ such that $A''(y_0)=0$. Then $y_0$ is a minimum of $A'(y)$ and a direct computation shows that $A'(y_0)>0$. We thus conclude that $A'(y)>0$ for all $y\in(-1,1)$.
\end{proof}

\end{document}
